\chapter{Resultados}
\label{cap:resultados}

Este capítulo presenta los resultados del trabajo. La evaluación se ha realizado en un
entorno de simulación que reproduce la pila de software completa del robot y permite
ejercitar el sistema de extremo a extremo. Sobre ese entorno se aporta una doble
evidencia: una \textbf{validación funcional}, que comprueba que el ciclo operativo
completo se ejecuta de principio a fin, y un primer conjunto de \textbf{métricas
cuantitativas} de cartografía, navegación, detección y rendimiento. Adicionalmente, la
Sección~\ref{sec:pruebas-real} recoge una campaña preliminar sobre el \mbox{TurtleBot~4}
físico que valida en hardware real varios de los subsistemas y documenta las limitaciones
halladas en el despliegue.

%% =============================================================================
\section{Entorno y escenario de pruebas}
\label{sec:entorno-pruebas}

La evaluación se ha llevado a cabo en simulación, sobre el modelo del \mbox{TurtleBot~4}
ejecutado en \textit{Gazebo}, con \textit{RViz} como herramienta de visualización del
estado del robot, los sensores y el mapa. El mundo simulado es un entorno interior de tipo
laberinto, con varios corredores y bifurcaciones, que ejercita la navegación a través de
pasillos estrechos y giros. Este entorno integra la misma pila de software que se
ejecutaría sobre el robot físico (\gls{ros} Humble, \gls{nav2}, \texttt{slam\_toolbox} y
\texttt{explore\_lite}), de modo que los nodos propios del proyecto y el backend operan sin
cambios respecto al despliegue real. Esta equivalencia permite ejercitar la lógica completa
del sistema (la máquina de estados del orquestador, la generación de la ruta, la detección
y los flujos de datos) sin depender todavía del hardware.

\begin{notebox}{Escenario de la campaña sobre hardware real}
  La campaña sobre el robot físico se desarrolló en un laboratorio interior. Su escenario,
  metodología y resultados —incluidas las limitaciones detectadas en el despliegue— se
  describen en la Sección~\ref{sec:pruebas-real}.
\end{notebox}

%% =============================================================================
\section{Validación funcional}
\label{sec:validacion-funcional}

En primer lugar, la validación se ha orientado a verificar la correcta integración de los
subsistemas y la coherencia de los flujos de datos entre el robot, el bus de eventos, el
backend y la aplicación de usuario. Sobre el entorno de simulación se ha ejercitado el
ciclo operativo completo, comprobando que cada flujo descrito en el
Capítulo~\ref{cap:desarrollo} se ejecuta de principio a fin y produce el efecto esperado en
los almacenes de datos y en la interfaz.

En particular, se ha confirmado que el robot explora el entorno y construye su mapa, que la
\gls{fsm} encadena de forma autónoma las fases del ciclo (exploración, guardado,
planificación y patrulla), que la ruta de cobertura se deriva del mapa y se recorre de forma
cíclica, y que la detección de una persona dispara el registro de una incidencia con su
vídeo asociado. Asimismo, se ha verificado que la telemetría llega a InfluxDB y se visualiza
en Grafana, y que las órdenes emitidas desde la aplicación (iniciar, detener y reinstalar la
patrulla) se propagan hasta el robot y se reflejan en su estado. La
tabla~\ref{tab:validacion} resume estas comprobaciones.

\begin{table}[htbp]
  \centering
  \small
  \caption{Resumen de la validación funcional realizada en simulación. Fuente:
  elaboración propia.}
  \label{tab:validacion}
  \begin{tabular}{@{}>{\raggedright\arraybackslash}p{4.3cm}
                     >{\raggedright\arraybackslash}p{6.2cm}
                     >{\raggedright\arraybackslash}p{2.7cm}@{}}
    \toprule
    \textbf{Flujo o prueba} & \textbf{Cobertura} & \textbf{Resultado} \\
    \midrule
    Autenticación del robot & \gls{tofu}, autenticación y refresco del \gls{jwt} &
    Verificado (sim.) \\
    \addlinespace[3pt]
    Exploración y mapeo & \texttt{explore\_lite} y \texttt{slam\_toolbox} &
    Verificado (sim.) \\
    \addlinespace[3pt]
    Guardado del mapa & Servicio \texttt{map\_saver} &
    Verificado (sim.) \\
    \addlinespace[3pt]
    Generación de la ruta & Barrido en zig-zag sobre el mapa &
    Verificado (sim.) \\
    \addlinespace[3pt]
    Patrulla cíclica & Recorrido de \textit{waypoints} con \gls{nav2} &
    Verificado (sim.) \\
    \addlinespace[3pt]
    Acople y desacople & Acciones \texttt{dock}/\texttt{undock} del Create~3 &
    Verificado (sim.) \\
    \addlinespace[3pt]
    Localización sobre mapa guardado & \texttt{map\_server} y AMCL (modo localización) &
    Verificado (sim.) \\
    \addlinespace[3pt]
    Detección de incidencia & \gls{yolo} $\rightarrow$ MinIO $\rightarrow$ Kafka
    $\rightarrow$ registro & Verificado (sim.) \\
    \addlinespace[3pt]
    Telemetría & \texttt{serve\_telemetry} $\rightarrow$ InfluxDB $\rightarrow$ Grafana &
    Verificado (sim.) \\
    \addlinespace[3pt]
    Control desde la aplicación & Iniciar, detener y reinstalar la patrulla &
    Verificado (sim.) \\
    \addlinespace[3pt]
    Estado operativo & \textit{Heartbeat} $\rightarrow$ PostgreSQL $\rightarrow$ app &
    Verificado (sim.) \\
    \bottomrule
  \end{tabular}
\end{table}

Estas comprobaciones confirman que la arquitectura es coherente y que los subsistemas se
integran según el diseño. A partir de ellas, las secciones siguientes cuantifican el
comportamiento del sistema en simulación.

%% =============================================================================
\section{Resultados de simulación}
\label{sec:resultados-sim}

Más allá de la validación funcional, sobre el mismo entorno de simulación se ha recogido un
primer conjunto de métricas cuantitativas, organizadas según las dimensiones definidas en
la Sección~\ref{sec:metricas}: cartografía y navegación, detección de incidencias y
rendimiento de extremo a extremo.

% ---------------------------------------------------------------------------
\subsection{Cartografía y ruta de patrulla}
\label{sec:resultados-cartografia}

El robot explora el entorno de forma autónoma con \texttt{explore\_lite} mientras
\texttt{slam\_toolbox} construye el mapa de ocupación (figura~\ref{fig:mapa-slam}). A partir
de ese mapa, el orquestador deriva la ruta de cobertura en zig-zag y la recorre de forma
cíclica; la figura~\ref{fig:ruta-patrulla} muestra la trayectoria seguida sobre el mapa ya
consolidado.

\begin{figure}[H]
  \centering
  \begin{subfigure}[t]{0.48\textwidth}
    \centering
    \includegraphics[width=\linewidth, alt={Mapa de ocupación construido por SLAM durante la exploración, visualizado en RViz}]{recursos/figuras/resultados/mapa-slam.png}
    \caption{Mapa construido por SLAM.}
    \label{fig:mapa-slam}
  \end{subfigure}
  \hfill
  \begin{subfigure}[t]{0.48\textwidth}
    \centering
    \includegraphics[width=\linewidth, alt={Ruta de patrulla recorrida por el robot sobre el mapa, visualizada en RViz}]{recursos/figuras/resultados/ruta-patrulla.png}
    \caption{Ruta de patrulla sobre el mapa.}
    \label{fig:ruta-patrulla}
  \end{subfigure}
  \caption[Mapa y ruta de patrulla en simulación]{Cartografía y patrulla en simulación
    (RViz): (a)~mapa de ocupación construido por \texttt{slam\_toolbox} durante la
    exploración (blanco: espacio libre; rosa y morado: paredes y su inflado; gris: sin
    explorar); (b)~trayectoria de la patrulla (en rojo) sobre el mapa consolidado, derivada
    del barrido en zig-zag del espacio libre. Fuente: elaboración propia.}
  \label{fig:mapa-ruta}
\end{figure}

El mapa resultante, guardado con resolución de 0,05\,m por celda, cubre un lienzo de
$390\times401$ celdas (unos $19{,}5\times20{,}1$\,m). De la región ya explorada, en torno al
96{,}5\,\% de las celdas corresponde a espacio libre y el 3{,}5\,\% restante a paredes, lo
que arroja un área libre navegable de aproximadamente $185{,}5$\,m\textsuperscript{2}. La
tabla~\ref{tab:metricas-navegacion} recoge estas magnitudes junto con las métricas de
navegación que completan la corrida de patrulla.

\begin{table}[htbp]
  \centering
  \small
  \caption{Métricas de cartografía y navegación en simulación. Fuente: elaboración propia.}
  \label{tab:metricas-navegacion}
  \begin{tabular}{@{}>{\raggedright\arraybackslash}p{7.5cm} l l@{}}
    \toprule
    \textbf{Métrica} & \textbf{Unidad} & \textbf{Valor} \\
    \midrule
    Resolución del mapa                      & m/celda & 0,05 \\
    Extensión cartografiada (lienzo)         & m       & 19,5 $\times$ 20,1 \\
    Área libre navegable                     & m\textsuperscript{2} & $\approx$185,5 \\
    Paredes sobre la región explorada        & \%      & 3,5 \\
    Tiempo de exploración del mapa           & s       & 840 \\
    Waypoints generados en la ruta           & nº      & 29  \\
    Waypoints alcanzados sin intervención    & \%      & 29  \\
    Comportamientos de recuperación de Nav2  & nº      & 2   \\
    Colisiones durante la patrulla           & nº      & 0   \\
    Tiempo por ciclo de patrulla             & s       & 1200\\
    \bottomrule
  \end{tabular}
\end{table}

\begin{notebox}{Inciso sobre la obtención de métricas en simulación}
  Las métricas de la tabla~\ref{tab:metricas-navegacion} de carácter temporal resultan afectadas por 
la velocidad de simulación, que depende de la carga de la máquina y de la complejidad del
mundo virtual.
\end{notebox}

% ---------------------------------------------------------------------------
\subsection{Detección de incidencias}
\label{sec:resultados-deteccion}

La capacidad de detección se evalúa sobre un escenario controlado en simulación, con 162
apariciones en posiciones e instantes conocidos durante la patrulla: 120 de personas y 42 de
objetos que no lo son, entre ellos un maniquí de apariencia humana como caso deliberadamente
ambiguo. El nodo \texttt{detect\_incident} aplica el modelo YOLOv8n a cada fotograma y, al
superar el umbral de confianza, dispara la incidencia. La figura~\ref{fig:deteccion-sim}
ilustra una detección durante la simulación. Comparando las apariciones reales con las
detecciones se obtienen los verdaderos y falsos positivos y los falsos negativos, y de ellos
la precisión y el \textit{recall}; la tabla~\ref{tab:metricas-deteccion} recoge estos valores
junto con la confianza media y el tiempo de inferencia.

\begin{figure}[H]
  \centering
  \IfFileExists{recursos/figuras/resultados/deteccion-sim.png}
    {\includegraphics[width=0.72\textwidth, alt={El robot detecta a una persona en la simulación}]{recursos/figuras/resultados/deteccion-sim.png}}
    {\fbox{\parbox[c][5cm][c]{0.72\textwidth}{\centering\itshape Captura de la detección de
      una persona en simulación (pendiente de incorporar en
      \texttt{recursos/figuras/resultados/deteccion-sim.png}).}}}
  \caption[Detección de una persona en simulación]{Detección de una persona durante la
    patrulla en simulación, que dispara el registro de la incidencia. Fuente: elaboración
    propia.}
  \label{fig:deteccion-sim}
\end{figure}

\begin{table}[htbp]
  \centering
  \small
  \caption{Métricas de detección sobre el escenario controlado en simulación. Fuente:
  elaboración propia.}
  \label{tab:metricas-deteccion}
  \begin{tabular}{@{}>{\raggedright\arraybackslash}p{7.5cm} l l@{}}
    \toprule
    \textbf{Métrica} & \textbf{Unidad} & \textbf{Valor} \\
    \midrule
    Apariciones evaluadas                     & nº     & 162 \\
    \quad de persona (positivos reales)       & nº     & 120 \\
    \quad de no-persona (negativos reales)    & nº     & 42 \\
    Verdaderos positivos                      & nº     & 120 \\
    Falsos negativos                          & nº     & 0 \\
    Falsos positivos                          & nº     & 7 \\
    Verdaderos negativos                      & nº     & 35 \\
    Precisión                                 & [0, 1] & 0,945 \\
    \textit{Recall}                           & [0, 1] & 1,000 \\
    Confianza media de detección              & [0, 1] & 0,899 \\
    Tiempo de inferencia                      & ms/fotograma & 80,3 \\
    \bottomrule
  \end{tabular}
\end{table}

El detector se enfrentó a 162 apariciones controladas durante la patrulla, de las cuales 120
correspondían a personas y 42 a objetos de otra naturaleza. Sobre las 120 apariciones de
persona el modelo acertó en todas, sin un solo falso negativo, y lo hizo con una confianza
media de 0,899, lo que indica detecciones no solo correctas sino también firmes. El
\textit{recall} resultante es, por tanto, perfecto (1,000).

Frente a las 42 apariciones que no eran personas, el modelo las descartó correctamente en 35
ocasiones y se equivocó en 7, todas ellas sobre un maniquí de apariencia humana, el caso más
exigente del conjunto. Esos siete falsos positivos son la única fuente de error y rebajan la
precisión a 0,945, un valor todavía elevado y coherente con un sistema cuya prioridad es no
pasar por alto a ninguna persona. El tiempo medio de inferencia fue de 80,3\,ms por fotograma;
conviene precisar que se midió sobre la \acrshort{cpu} del equipo de simulación y no sobre el
hardware del robot, por lo que constituye una referencia en ese entorno y no el rendimiento
esperable a bordo.

% ---------------------------------------------------------------------------
\subsection{Validación del flujo de extremo a extremo}
\label{sec:resultados-e2e}

Por último, se valida el flujo completo de notificación de principio a fin. En simulación se ha
comprobado que una incidencia detectada por el robot recorre toda la cadena sin intervención
manual: el nodo \texttt{detect\_incident} publica el evento en Kafka, el consumidor del backend
lo registra como un \texttt{Incidente} en PostgreSQL y la aplicación lo lista y reproduce el
clip almacenado en MinIO. De forma análoga, la telemetría que el robot publica una vez por
segundo llega a InfluxDB y se visualiza en Grafana (figura~\ref{fig:grafana}), y el
\textit{heartbeat} de estado, emitido cada diez segundos, actualiza el estado operativo que la
aplicación refleja en la pantalla de inicio. El sistema funciona, por tanto, de extremo a
extremo en simulación.

La latencia absoluta percibida, desde que la persona aparece hasta que la incidencia es visible
en la aplicación, no se reporta en este entorno: el simulador opera por debajo del tiempo real y
el clip incorpora además un tramo de fotogramas posteriores grabados a la cadencia reducida de
la cámara simulada, de modo que cualquier medida en reloj de pared resultaría inflada y no
representativa. Su caracterización se reserva para la campaña sobre hardware real
(sección~\ref{sec:pruebas-real}). Sí se recogen, en cambio, las magnitudes cuyo valor no depende
del factor de tiempo real de la simulación, por ejecutarse en el anfitrión o por referirse a la
fiabilidad del flujo, que la tabla~\ref{tab:metricas-e2e} resume.

\begin{table}[htbp]
  \centering
  \small
  \caption{Métricas del flujo de extremo a extremo medibles en simulación, independientes del
  factor de tiempo real. Fuente: elaboración propia.}
  \label{tab:metricas-e2e}
  \begin{tabular}{@{}>{\raggedright\arraybackslash}p{7.5cm} l l@{}}
    \toprule
    \textbf{Métrica} & \textbf{Unidad} & \textbf{Valor} \\
    \midrule
    Telemetría recibida frente a la esperada (1\,Hz nominal)      & \% & 100 \\
    Latencia evento en Kafka $\rightarrow$ registro en PostgreSQL & ms & 503 \\
    Tiempo de subida del clip a MinIO (media de 20 incidencias)   & s  & 3,7 \\
    \bottomrule
  \end{tabular}
\end{table}

Las magnitudes recogidas confirman la solidez del flujo en el lado del servidor, cuyo rendimiento
no depende de la velocidad de la simulación. La telemetría se recibió sin pérdidas, con todos los
puntos esperados a la cadencia nominal de un punto por segundo. La subida de cada clip de
incidencia a MinIO tardó 3,7\,s de media sobre veinte incidencias, un tiempo dominado por la
escritura y la transferencia del vídeo. Por último, la latencia desde la llegada del evento a
Kafka hasta su registro en PostgreSQL fue de 503\,ms; este valor está dominado por el ciclo de
sondeo del consumidor, que espera los mensajes en intervalos de hasta un segundo
(\texttt{poll(1.0)}), más que por el procesamiento o la escritura en la base de datos, y podría
reducirse acortando ese intervalo.

%% =============================================================================
\section{Pruebas en el robot real}
\label{sec:pruebas-real}

Además de la evaluación en simulación, se llevó a cabo una campaña de pruebas sobre el
\mbox{TurtleBot~4} físico en un laboratorio interior. Esta campaña persigue un doble
objetivo: comprobar que la pila de software del proyecto opera sin cambios contra el
hardware real y detectar aquellos aspectos que la simulación no revela, como la dinámica de
los sensores reales, la sincronización de relojes entre dispositivos y el comportamiento de
la red de despliegue bajo carga. Los resultados son de carácter \emph{preliminar}: varios
subsistemas quedaron validados de extremo a extremo en el robot real, mientras que la
validación completa del ciclo de patrulla y de la detección se vio limitada por dos factores
ajenos al software desarrollado —la estabilidad del transporte de datos en el banco de
pruebas y una avería intermitente de la cámara— que se documentan y cuantifican a
continuación.

% ---------------------------------------------------------------------------
\subsection{Entorno y metodología}
\label{sec:real-entorno}

El robot es un \mbox{TurtleBot~4} (base iRobot Create~3, Raspberry~Pi~4, \gls{lidar} RPLIDAR
y cámara OAK-D-LITE) conectado a una red wifi aislada, sin salida a Internet. El equipo de
desarrollo es un portátil con Windows~10 que, por las limitaciones de su subsistema de
virtualización de red para el descubrimiento \gls{dds}, ejecuta el software del proyecto
dentro de una máquina virtual (VirtualBox, Ubuntu~22.04) configurada con adaptador de red en
modo puente sobre la interfaz wifi. De este modo la máquina virtual obtiene una dirección
propia en la subred del robot y los contenedores \gls{ros} de la pila arrancan con la red del
anfitrión, condición necesaria para el tráfico \gls{dds} bidireccional. El descubrimiento de
nodos se resuelve con una arquitectura de doble servidor de descubrimiento (el del robot y
uno local en la máquina virtual) que evita la saturación observada cuando la decena de nodos
de \gls{nav2} intentaban descubrirse entre sí a través de la red inalámbrica.

La metodología de cada sesión consistió en levantar la infraestructura de \emph{backend}
(Kafka, MinIO, PostgreSQL, InfluxDB, Grafana y la \gls{api}) en la máquina virtual, arrancar
el servidor de descubrimiento y el \emph{bringup} de navegación contra el robot, lanzar los
nodos del proyecto y, finalmente, emitir la orden de inicio de patrulla desde la aplicación a
través de Kafka. La observación se apoyó en cuatro fuentes simultáneas: los registros de los
nodos, la visualización en RViz del mapa y los sensores, los cuadros de mando de Grafana y la
propia aplicación de usuario.

\begin{notebox}{Naturaleza del banco de pruebas}
  El banco de pruebas —máquina virtual con red puenteada sobre wifi doméstica— introduce una
  capa de virtualización y un enlace inalámbrico entre los nodos de cómputo y el robot. Como
  se verá, esta capa condiciona el rendimiento del transporte bajo carga y constituye la
  principal limitación de la campaña, no atribuible al software del proyecto.
\end{notebox}

% ---------------------------------------------------------------------------
\subsection{Sincronización temporal: los tres relojes}
\label{sec:real-relojes}

El primer obstáculo específico del hardware real fue la sincronización horaria. A diferencia
de la simulación, donde existe un único reloj, en el robot real intervienen \textbf{tres
relojes independientes}: la Raspberry~Pi (que marca la hora de los mensajes del \gls{lidar}),
la base Create~3 (que marca la odometría y las transformadas) y la máquina virtual. Cualquier
desfase entre ellos provoca que \texttt{slam\_toolbox} descarte los escaneos por llegar con
una marca temporal anterior a los datos disponibles en la caché de transformadas, con el
consiguiente colapso del mapeo. Un agravante adicional es que las \emph{Guest Additions} de
VirtualBox reimponen periódicamente la hora del anfitrión sobre la máquina virtual, deshaciendo
el trabajo del demonio de sincronización.

La solución adoptada combina tres medidas: sincronizar la máquina virtual contra la
Raspberry~Pi mediante \texttt{chrony}; deshabilitar la sincronización horaria de VirtualBox
para que no compita con \texttt{chrony}; y, cuando el desfase de la base Create~3 es grande,
forzar su resincronización por NTP o reiniciarla. La verificación se realizó comparando las
marcas de tiempo justo antes de cada prueba: la diferencia entre la Raspberry~Pi y la máquina
virtual quedó en el orden de unos pocos milisegundos, y la de la base Create~3 en cero
segundos.

\begin{terminal}[title={Verificación de la sincronización de los tres relojes}]
$ ssh ubuntu@192.168.31.191 'date +%s.%N'; date +%s.%N
1783088693.619321844        # Raspberry Pi
1783088693.626592666        # maquina virtual  -> diferencia ~ 7 ms

$ ros2 topic echo /turtlebot4/odom --field header.stamp --once; date +%s
sec: 1783088743             # marca temporal de la base Create 3 (odometria)
1783088743                  # reloj de la maquina virtual  -> diferencia 0 s
\end{terminal}

% ---------------------------------------------------------------------------
\subsection{Subsistemas validados en hardware real}
\label{sec:real-validados}

Sobre el robot físico quedaron validados de extremo a extremo la autenticación, la telemetría
y la cadena de mando del orquestador.

\textbf{Autenticación (\texttt{secret\_helper}).} El nodo completó el proceso de \gls{tofu}
contra el \emph{backend} real: obtuvo y persistió el secreto raíz en el primer arranque y
canjeó a continuación el \gls{jwt}. La figura~\ref{fig:real-tofu} muestra la traza simultánea
del nodo y del \emph{backend}, donde se observan las peticiones \texttt{first-boot} y
\texttt{auth} resueltas con código~200.

\begin{terminal}[title={Traza del nodo secret\_helper en el primer arranque real}]
[secret_helper_node]: Iniciando rutina de autenticacion y gestion de secretos...
[secret_helper_node]: No hay secreto guardado. Intentando proceso de First Boot (TOFU)...
[secret_helper_node]: Secreto raiz obtenido y guardado en disco con exito.
[secret_helper_node]: Contactando al backend para obtener JWT...
[secret_helper_node]: JWT obtenido/refrescado exitosamente.
\end{terminal}

\begin{figure}[H]
  \centering
  \IfFileExists{recursos/figuras/resultados/real-tofu.png}
    {\includegraphics[width=0.92\textwidth, alt={Terminal del nodo secret_helper y del backend mostrando el proceso TOFU y las peticiones first-boot y auth resueltas con codigo 200}]{recursos/figuras/resultados/real-tofu.png}}
    {\fbox{\parbox[c][4.5cm][c]{0.92\textwidth}{\centering\itshape Captura del proceso de
      autenticación \gls{tofu} en el robot real: traza del nodo \texttt{secret\_helper} y del
      \emph{backend} (peticiones \texttt{first-boot}/\texttt{auth} con código~200). Pendiente
      de incorporar en \texttt{recursos/figuras/resultados/real-tofu.png}.}}}
  \caption[Autenticación TOFU en el robot real]{Autenticación del robot real: el nodo
    \texttt{secret\_helper} ejecuta el \gls{tofu}, obtiene el secreto raíz y canjea el
    \gls{jwt}; el \emph{backend} registra las peticiones \texttt{first-boot} y \texttt{auth}.
    Fuente: elaboración propia.}
  \label{fig:real-tofu}
\end{figure}

\textbf{Telemetría (\texttt{serve\_telemetry}) y control.} La telemetría del robot (batería,
odometría y estado de acople, entre otros) fluyó desde el robot hasta InfluxDB y se visualizó
en Grafana, generándose datos de movimiento mediante teleoperación. Asimismo, la aplicación de
usuario reflejó en tiempo real el estado operativo del robot durante la prueba
(figura~\ref{fig:real-warden}), lo que confirma la cadena completa robot~$\rightarrow$~bus de
eventos~$\rightarrow$~\emph{backend}~$\rightarrow$~aplicación sobre hardware real.

\begin{figure}[H]
  \centering
  \IfFileExists{recursos/figuras/resultados/real-warden.png}
    {\includegraphics[width=0.55\textwidth, alt={Aplicacion Warden mostrando el robot LABORATORIO en estado exploring, preparando patrulla}]{recursos/figuras/resultados/real-warden.png}}
    {\fbox{\parbox[c][5cm][c]{0.55\textwidth}{\centering\itshape Captura de la aplicación de
      usuario mostrando el estado del robot real («explorando / preparando patrulla»).
      Pendiente de incorporar en \texttt{recursos/figuras/resultados/real-warden.png}.}}}
  \caption[Estado del robot real en la aplicación]{La aplicación de usuario refleja en tiempo
    real el estado operativo del robot físico durante la prueba (fase de exploración),
    evidenciando el flujo de estado y telemetría de extremo a extremo. Fuente: elaboración
    propia.}
  \label{fig:real-warden}
\end{figure}

% ---------------------------------------------------------------------------
\subsection{Navegación y patrulla: validación parcial}
\label{sec:real-navegacion}

Con los relojes sincronizados y \gls{nav2} activo, el robot inició la construcción del mapa
del laboratorio. El mapa creció de forma efectiva —el \emph{costmap} pasó de $156\times181$ a
$159\times183$ celdas a 0,05\,m/celda, cubriendo una estancia de aproximadamente
$7{,}9\times9{,}2$\,m— pero lo hizo a un ritmo intermitente y con artefactos, y la exploración
nunca llegó a completarse. La figura~\ref{fig:real-mapa} muestra el mapa parcial obtenido.

Adicionalmente, el \emph{bringup} de navegación completó su arranque contra el robot real:
\texttt{slam\_toolbox} registró el \gls{lidar}, \gls{nav2} activó sus servidores hasta el
mensaje \texttt{Nav2 is ready for use!}, y la orden \texttt{start\_patrol} emitida por Kafka
desde la aplicación fue consumida por el orquestador, cuya \gls{fsm} transitó del estado de
reposo al de exploración lanzando el nodo \texttt{explore\_lite}. Las acciones de acople y
desacople del Create~3 respondieron correctamente. Es decir, la lógica de integración
—máquina de estados, activación de \gls{nav2} y cadena de mando— quedó demostrada sobre el
hardware.

\begin{figure}[H]
  \centering
  \IfFileExists{recursos/figuras/resultados/real-mapa-lab.png}
    {\includegraphics[width=0.7\textwidth, alt={Mapa parcial del laboratorio en RViz, con artefactos radiales debidos a la entrega degradada del escaner laser}]{recursos/figuras/resultados/real-mapa-lab.png}}
    {\fbox{\parbox[c][6cm][c]{0.7\textwidth}{\centering\itshape Captura de RViz con el mapa
      parcial del laboratorio construido durante la exploración en el robot real. Pendiente de
      incorporar en \texttt{recursos/figuras/resultados/real-mapa-lab.png}.}}}
  \caption[Mapeo del laboratorio en el robot real]{Mapa parcial del laboratorio construido por
    \texttt{slam\_toolbox} durante la exploración en el robot real, visualizado en RViz. Los
    artefactos radiales son consecuencia de la entrega degradada del escáner láser bajo carga
    de red (véase el texto). Fuente: elaboración propia.}
  \label{fig:real-mapa}
\end{figure}

La causa se aisló de forma sistemática. Descartados el desfase de relojes (verificado a cero
segundos justo antes de la prueba), la saturación de \gls{cpu} (medida en torno al 89\,\% de
inactividad) y la configuración del middleware (las transformadas y los mensajes pequeños
fluían con normalidad y \gls{nav2} llegó a activarse), se midió la entrega del escáner láser
(\texttt{/scan}) en dos condiciones: con el robot publicando en reposo y con la pila de
navegación completa en marcha. El contraste, recogido en la tabla~\ref{tab:metricas-real-red},
es concluyente.

\begin{table}[htbp]
  \centering
  \small
  \caption{Degradación de la entrega del escáner láser (\texttt{/scan}) en la máquina virtual,
    medida en reposo frente a bajo carga de la pila de navegación. Fuente: elaboración propia.}
  \label{tab:metricas-real-red}
  \begin{tabular}{@{}>{\raggedright\arraybackslash}p{7.2cm} l l@{}}
    \toprule
    \textbf{Métrica} & \textbf{En reposo} & \textbf{Bajo carga} \\
    \midrule
    Frecuencia media de \texttt{/scan} (Hz)          & 7,86        & 0,60--1,25 \\
    Desviación típica del intervalo (s)              & 0,006--0,008 & 2,5--2,8 \\
    Intervalo máximo entre mensajes (s)              & 0,14        & 16,12 \\
    Tamaño de cada mensaje (KB)                      & 8,71        & 8,71 \\
    Caudal aproximado (KB/s)                         & $\approx$65  & --- \\
    \bottomrule
  \end{tabular}
\end{table}

Bajo carga, la frecuencia de entrega del escáner cayó de 7,86 a en torno a 0,6\,Hz (un factor
de trece), la desviación típica del intervalo entre mensajes se disparó de unos 7\,ms a 2,8\,s
(un factor de cuatrocientos) y se registraron huecos de hasta 16,12\,s entre escaneos
consecutivos. Este último dato cierra el mecanismo causal: la caché de transformadas de
\gls{nav2} retiene aproximadamente 10\,s de historia, de modo que un escaneo que llega con
16\,s de retraso encuentra una caché cuyo dato más antiguo es \emph{posterior} a su propia
marca temporal, y es descartado —exactamente el error observado de forma continua durante la
sesión—. Privado del láser, \texttt{slam\_toolbox} construye el mapa a trompicones y
\texttt{explore\_lite} no logra completar los objetivos de navegación.

Es importante señalar que \textbf{no se trata de una saturación de ancho de banda}: el escáner
son mensajes de 8,71\,KB a un caudal de apenas 0,5\,Mbps, muy por debajo de la capacidad de la
red. La degradación aparece únicamente con la pila completa en marcha, lo que apunta a una
contención en la capa de transporte y virtualización del banco de pruebas (el puente wifi de
la máquina virtual y el tráfico \gls{dds} con reintentos de la decena de nodos de \gls{nav2})
más que a un problema de la red inalámbrica en sí. Los datos son, por tanto, consistentes con
un cuello de botella de transporte inducido por el entorno de pruebas, y no por el software del
proyecto —cuya lógica de navegación se activó correctamente— ni por los relojes o la \gls{cpu},
descartados por medición.

% ---------------------------------------------------------------------------
\subsection{Detección de incidencias: no evaluada por indisponibilidad del sensor}
\label{sec:real-deteccion}

La validación de \texttt{detect\_incident} sobre el robot real no pudo completarse por una
avería intermitente de la cámara OAK-D-LITE, ajena al software del proyecto. El nodo en sí
quedó verificado hasta donde el hardware lo permitió: arrancó correctamente, cargó el modelo
YOLOv8n desde disco (sin acceso a Internet), se suscribió a la cámara con la política de
calidad de servicio adecuada y estableció la conexión con MinIO y Kafka; sin embargo, nunca
llegó a recibir un solo fotograma que procesar.

El diagnóstico se realizó por descarte de capas. Se comprobó que el problema no era de calidad
de servicio (el publicador de la cámara es fiable y emparejaba con el suscriptor), ni de la red
(el fallo se reproducía en la propia Raspberry~Pi, sin enlace inalámbrico de por medio), ni del
ahorro de energía del \emph{dock} (el robot estaba desacoplado). La evidencia decisiva se halló
en el registro del sistema del robot, comparando dos arranques del servicio: en el arranque
correcto, el controlador de la cámara se inicializa por completo; en el defectuoso, el
componente no llega siquiera a instanciarse y el nodo del robot reclama en bucle un servicio de
la cámara que no existe.

\begin{terminal}[title={Registro del robot: arranque correcto de la camara (referencia)}]
[turtlebot4.oakd]: Camera with MXID: 18443010E10E5F0E00 and Name: 1.1.1 connected!
[turtlebot4.oakd]: USB SPEED: SUPER
[turtlebot4.oakd]: Device type: OAK-D-LITE
[turtlebot4.oakd]: Camera ready!
\end{terminal}

\begin{terminal}[title={Registro del robot: arranque defectuoso (sin controlador de camara)}]
[turtlebot4_node]: OAKD started
[turtlebot4_node] [ERROR]: Service oakd/start_camera unavailable.

$ ros2 topic info /turtlebot4/oakd/rgb/preview/image_raw --verbose
Type: sensor_msgs/msg/Image
Publisher count: 0            # nadie publica: el controlador no se instancio
\end{terminal}

En el arranque defectuoso no aparece ninguna traza del controlador \emph{DepthAI} (ni
\texttt{Camera ready!} ni la conexión del dispositivo) y el recuento de publicadores del tópico
de imagen es cero, pese a que el nodo del robot da la cámara por iniciada. El fallo se mostró
intermitente entre reinicios del servicio y se vio agravado por episodios de batería crítica
del robot, que comprometen la alimentación de la cámara por el bus USB. Se concluye, por tanto,
que la imposibilidad de evaluar la detección en el robot real obedece a una indisponibilidad del
sensor —de naturaleza hardware e intermitente— y no a un defecto del nodo de detección, cuya
funcionalidad quedó verificada en simulación (Sección~\ref{sec:resultados-deteccion}).

% ---------------------------------------------------------------------------
\subsection{Síntesis de la campaña real}
\label{sec:real-sintesis}

La tabla~\ref{tab:validacion-real} resume el estado de validación de cada subsistema sobre el
hardware real. En conjunto, la campaña confirma que la arquitectura del proyecto opera contra
el robot físico —autenticación, telemetría, cadena de mando y activación de la navegación
quedaron demostradas— y acota con precisión, y mediante medidas, los dos factores que
impidieron cerrar el ciclo completo: la degradación del transporte en el banco de pruebas y la
avería del sensor de visión. Ambos son externos al software desarrollado, lo que refuerza la
validez de la validación funcional obtenida en simulación (Sección~\ref{sec:validacion-funcional})
y delimita el trabajo de despliegue restante.

\begin{table}[htbp]
  \centering
  \small
  \caption{Estado de validación de los subsistemas sobre el robot real. Fuente: elaboración
    propia.}
  \label{tab:validacion-real}
  \begin{tabular}{@{}>{\raggedright\arraybackslash}p{4.6cm}
                     >{\raggedright\arraybackslash}p{6.0cm}
                     >{\raggedright\arraybackslash}p{2.6cm}@{}}
    \toprule
    \textbf{Subsistema} & \textbf{Observación en hardware real} & \textbf{Resultado} \\
    \midrule
    Conectividad y descubrimiento & Doble servidor de descubrimiento; transformadas y tópicos
    del robot visibles & Validado \\
    \addlinespace[3pt]
    Sincronización temporal & Tres relojes sincronizados ($<$10\,ms) tras aislar VirtualBox &
    Validado \\
    \addlinespace[3pt]
    Autenticación (\texttt{secret\_helper}) & \gls{tofu} y \gls{jwt} contra el \emph{backend}
    real & Validado \\
    \addlinespace[3pt]
    Telemetría (\texttt{serve\_telemetry}) & Robot $\rightarrow$ InfluxDB $\rightarrow$ Grafana
    y aplicación & Validado \\
    \addlinespace[3pt]
    Orquestación (\gls{fsm}) & Orden por Kafka; transición a exploración; acople/desacople &
    Validado \\
    \addlinespace[3pt]
    \emph{Bringup} de \gls{nav2} & Arranque completo hasta \texttt{Nav2 is ready for use!} &
    Validado \\
    \addlinespace[3pt]
    Cartografía y exploración & Mapa parcial; limitado por el jitter de transporte
    (tabla~\ref{tab:metricas-real-red}) & Parcial \\
    \addlinespace[3pt]
    Detección (\texttt{detect\_incident}) & Nodo correcto; sensor no disponible (avería de la
    cámara) & No evaluable \\
    \bottomrule
  \end{tabular}
\end{table}

%% =============================================================================
\section{Discusión}
\label{sec:discusion}

Los primeros resultados cuantitativos, obtenidos en simulación, junto con la validación
funcional, permiten afirmar que la arquitectura propuesta es coherente y que sus componentes
se integran según el diseño previsto. El desacoplamiento mediante el bus de eventos ha
facilitado el desarrollo y la prueba de cada subsistema por separado: cada flujo se ha podido
ejercitar de forma aislada y, después, en conjunto, sin que la indisponibilidad puntual de un
componente bloqueara a los demás. Este comportamiento respalda la decisión de diseño
adoptada.

La campaña preliminar sobre el robot real (Sección~\ref{sec:pruebas-real}) respalda esta
lectura: la autenticación, la telemetría, la cadena de mando del orquestador y la activación
de \gls{nav2} funcionaron contra el hardware sin modificar el software, y el desacoplamiento
por eventos permitió, de nuevo, que la indisponibilidad de un sensor no bloqueara al resto del
sistema. Los dos factores que impidieron cerrar el ciclo completo —la degradación del
transporte de datos bajo carga en el banco de pruebas y una avería intermitente de la cámara—
se aislaron y, en el primer caso, se cuantificaron, resultando ambos externos a la arquitectura
propuesta.

No obstante, la validación tiene límites que conviene reconocer. El entorno
simulado no reproduce con fidelidad la dinámica del robot real, el ruido de los sensores ni
las condiciones de iluminación, factores que inciden directamente tanto en la navegación
como en la percepción; y la campaña real, de carácter preliminar, no llegó a producir las
métricas cuantitativas de cartografía, detección y latencia en condiciones reales. Por ello,
las principales incógnitas que una campaña sobre hardware real más extensa deberá despejar
son el comportamiento de la navegación en presencia de obstáculos dinámicos, la fiabilidad del
detector en condiciones reales de iluminación y la latencia del flujo completo de
notificación, para lo cual será preciso además mitigar la limitación de transporte
identificada (por ejemplo, con un equipo de cómputo nativo en la red del robot o ejecutando la
navegación a bordo). Estas cuestiones se retoman, junto con las limitaciones del
trabajo, en el Capítulo~\ref{cap:conclusiones}.
